\section{Polynomial reduction}\label{sec:polynomial-reduction}

\subsection{One-step reduction}

Continue to use the fixed monomial order $m$.

\begin{definition}[One-step polynomial reduction]
\label{def:one-step-reduction}
Let $p,f,g\in\Rpoly$.  We say that $f$ reduces to $g$ by $p$ if
there are monomials $x^t$ and $x^s$ and a coefficient $c\in R$ such that
\begin{enumerate}[label=(\roman*),leftmargin=2em]
  \item $p\neq0$ and $x^t$ occurs in $f$;
  \item $\LM(p)\mid x^t$; equivalently, $x^t=x^s\LM(p)$;
  \item $c\LC(p)$ equals the coefficient of $x^t$ in $f$; and
  \item $g=f-cx^s p$.
\end{enumerate}
Reduction modulo a set $G$ means reduction by some polynomial $p\in G$.
The Lean definition records the exponents $s$ and $t$ explicitly and writes
the coefficient of $x^t$ in $f$ as \lean{f.coeff t}.
\end{definition}

The exact Lean definition is shown in
Listing~\ref{lst:reduction-definition}.

\begin{lstlisting}[style=lean,caption={One-step polynomial reduction and its existential closures.},label={lst:reduction-definition}]
def ReducesToBy
    (p f g : MvPolynomial σ R) (t : σ →₀ ℕ) : Prop :=
  p ≠ 0 ∧
    t ∈ f.support ∧
      ∃ s : σ →₀ ℕ,
        s + m.degree p = t ∧
          ∃ c : R,
            c * m.leadingCoeff p = f.coeff t ∧
              g = f - monomial s c * p

def ReducesToPoly (p f g : MvPolynomial σ R) : Prop :=
  ∃ t : σ →₀ ℕ, m.ReducesToBy p f g t

def ReducesToSet
    (P : Set (MvPolynomial σ R))
    (f g : MvPolynomial σ R) : Prop :=
  ∃ p ∈ P, m.ReducesToPoly p f g
\end{lstlisting}

\begin{remark}[Single-reducer strong reduction]
\label{rem:strong-reduction}
The relation in Definition~\ref{def:one-step-reduction} is a single-reducer
strong reduction relation: one step subtracts a monomial multiple of one
polynomial $p\in G$.  It does not implement weak reduction over coefficient
rings, where leading coefficients of several reducers may be combined to
cancel a single term.  For example, in $\mathbb{Z}[x]$, let
\[
  G=\{2x,3x\},\qquad f=x.
\]
Under the present relation, $x$ is normal modulo $G$, because neither $2$ nor
$3$ divides its coefficient $1$.  In a weak reduction framework, however, the
identity
\[
  x=(-1)(2x)+(3x)
\]
allows the term to be cancelled by combining the two reducers.
\end{remark}

\begin{definition}[Reduction closures]
\label{def:reduction-closures}
For a fixed set $G\subseteq\Rpoly$:
\begin{enumerate}[label=(\roman*),leftmargin=2em]
  \item $f\red{G}g$ denotes one-step reduction by some $p\in G$;
  \item $f\redstar{G}g$ denotes the reflexive-transitive closure of
        one-step reduction;
  \item $f\eqvred{G}g$ denotes the smallest equivalence relation containing
        one-step reduction; and
  \item $f\joinred{G}g$ means that there exists $h$ with
        $f\redstar{G}h$ and $g\redstar{G}h$.
\end{enumerate}
\end{definition}

The relation $\red{G}$ permits cancellation of any reducible monomial.
Mathlib's division procedure, by contrast, selects terms according to a
specific algorithm.  The arbitrary-term relation is used for confluence and
Church--Rosser statements; leading-term reduction remains available as the
separate predicate \lean{LTermReducesToPoly}.

\subsection{Reducibility over rings}

\begin{proposition}[Reducibility criterion over a commutative ring]
\label{prop:reducibility-ring}
For $p,f\in\Rpoly$, the polynomial $f$ is reducible by $p$ if and only if
\[
  p\neq0
  \quad\text{and}\quad
  \exists t\in\supp(f),\qquad
    \LM(p)\mid x^t
    \quad\text{and}\quad
    \LC(p)\mid\operatorname{coeff}_t(f).
\]
Since $\LM(p)=x^{\deg_m(p)}$, the monomial-divisibility condition is
equivalently the componentwise inequality $\deg_m(p)\leq t$ on exponent
vectors.  Thus reducibility over a general coefficient ring consists of both
monomial divisibility and coefficient divisibility.
\end{proposition}

Proposition~\ref{prop:reducibility-ring} is the exact content of
\lean{MonomialOrder.ReducesToPoly.reducible_iff_exists_degree_le_and_leadingCoeff_dvd}.
In particular, a unit leading coefficient is not required to define a
reduction step.  It is a sufficient condition that makes the coefficient
divisibility requirement automatic and reduces the criterion to monomial
divisibility alone.

\begin{remark}[The unit-leading-coefficient hypothesis]
\label{rem:unit-condition}
Several later translation and confluence theorems assume
\begin{equation}\label{eq:unit-condition}
  \forall g\in G,\qquad
  g\neq0\;\Longrightarrow\;\LC(g)\text{ is a unit}.
\end{equation}
The Lean statements encode the same condition as
\(
  \operatorname{IsUnit}(\LC(g))\lor g=0
\), which avoids imposing a unit condition on \(\LC(0)=0\).  Under this
hypothesis, whenever $\LM(g)$ divides a selected monomial, the required
coefficient multiple exists.  Over a field, Condition~\eqref{eq:unit-condition}
is automatic.  Over a general ring, it is an algebraic hypothesis used by
specific theorems rather than by the definition of reduction itself.
\end{remark}

\subsection{Termination by support colex order}

Total degree is not a termination measure for arbitrary-term reduction:
cancelling one monomial can introduce several monomials of the same total
degree.

\begin{definition}[Support colexicographic order]
\label{def:support-colex}
Identify the support of a polynomial with the finite set of monomials that
occur with nonzero coefficient.  If
$\supp(f)\mathbin{\triangle}\supp(g)\neq\varnothing$, let
\[
  M=\max\nolimits_m\bigl(\supp(f)\mathbin{\triangle}\supp(g)\bigr)
\]
be the largest monomial in the symmetric difference with respect to the fixed
monomial order.  We declare
\[
  f<_{\mathrm{supp}}g
\]
when $M\in\supp(g)$.  If the supports are equal, their symmetric difference is
empty and neither polynomial is smaller than the other.  Thus the polynomial
containing the largest monomial at which the two supports differ is the larger
one.
\end{definition}

In Lean, monomials are stored as exponent vectors.  The implementation
transports these exponents to the linearly ordered synonym attached to $m$
and then applies the finite-set colexicographic order; this is the role of
\lean{m.toSyn} and \lean{toColex} in Listing~\ref{lst:termination}.

\begin{lstlisting}[style=lean,caption={The support-colex measure and termination of polynomial reduction.},label={lst:termination}]
def supportColexLt (f g : MvPolynomial σ R) : Prop :=
  letI : LinearOrder m.syn := m.linearOrderSyn
  toColex (f.support.image m.toSyn) <
    toColex (g.support.image m.toSyn)

theorem supportColexLt_wellFounded :
    WellFounded
      (m.supportColexLt :
        MvPolynomial σ R → MvPolynomial σ R → Prop)

theorem ReducesToSet.wellFounded
    (P : Set (MvPolynomial σ R)) :
    WellFounded (flip (m.ReducesToSet P)) :=
  m.supportColexLt_wellFounded.mono
    (by
      intro f g h
      exact ReducesToSet.supportColexLt m P h)
\end{lstlisting}

Suppose that $f\red{G}g$ cancels the monomial $x^t$.  Then
$x^t\in\supp(f)\setminus\supp(g)$, while every monomial that occurs in $g$
but not in $f$ is strictly smaller than $x^t$.  Hence $x^t$ is the largest
monomial at which the supports differ, and it belongs to $f$.  Therefore
$g<_{\mathrm{supp}}f$.

\begin{theorem}[Termination of polynomial reduction]\label{thm:reduction-termination}
For every set $G\subseteq\Rpoly$, there is no infinite forward reduction
chain
\[
  f_0\red{G}f_1\red{G}f_2\red{G}\cdots.
\]
\end{theorem}

Listing~\ref{lst:termination} expresses this theorem in Lean as
well-foundedness of the flipped one-step reduction relation.

\subsection{Translation to ideal congruence}

If $f\red{G}g$, then the definition gives
$f-g=cx^s p$ for some $p\in G$; hence $f-g\in\ideal{G}$.  The same conclusion
holds for a finite reduction sequence and for the equivalence closure of
reduction.

Conversely, assume Condition~\eqref{eq:unit-condition} and
$f-g\in\ideal{G}$.  Here \emph{span induction} means proving a property for
every element of $\ideal{G}$ by checking it first for each generator in $G$
and then showing that it is preserved by zero, addition, and polynomial
scalar multiplication.  For a polynomial $x$, let $Q(x)$ be the assertion
that
\[
  a+hx\eqvred{G}a
\]
for all polynomials $a$ and $h$.  If $p\in G$, then
$hp\redstar{G}0$; hence
\[
  (a+hp)-a=hp\redstar{G}0,
\]
and the difference-to-equivalence part of the polynomial translation lemma,
applied to $(a+hp)-a$, gives $a+hp\eqvred{G}a$.  In Lean, this implication is
\lean{MonomialOrder.eqvGen_of_sub_reflTransGen_zero}; the generator case is
packaged as \lean{MonomialOrder.eqvGen_add_mul_mem}.  Thus $Q$ holds for the
generators and is preserved by zero, addition, and polynomial scalar
multiplication.  Span induction therefore gives $Q(f-g)$, and the resulting
full equivalence is
\leandoc{Mathlib/RingTheory/MvPolynomial/PolynomialReductions}{MonomialOrder.eqvGen_iff_sub_mem_span}{MonomialOrder.eqvGen_iff_sub_mem_span}.
Taking
$a=g$ and $h=1$ yields
\[
  f=g+(f-g)\eqvred{G}g.
\]

\begin{lstlisting}[style=lean,caption={Reduction equivalence coincides with congruence modulo the generated ideal.},label={lst:translation}]
theorem eqvGen_iff_sub_mem_span
    (P : Set (MvPolynomial σ R))
    (hP₀ :
      ∀ p ∈ P,
        IsUnit (m.leadingCoeff p) ∨ p = 0)
    {f g : MvPolynomial σ R} :
    f ⟷*[m, P] g ↔ f - g ∈ Ideal.span P
\end{lstlisting}

\begin{theorem}[Reduction/ideal translation]\label{thm:translation}
Assume \eqref{eq:unit-condition}.  Then
\[
  f\eqvred{G}g
  \quad\Longleftrightarrow\quad
  f-g\in\ideal{G}.
\]
\end{theorem}
